\chapter{Experimental Setup}
\label{chap:setup}

\section{Evaluation questions}

The experiments address three questions. First, does the fixed region graph
improve classification relative to the shared Swin-T baseline? Second, does
the local prototype soft assignment improve on the fixed grouping under the
same task protocol? Third, are any observed benefits consistent across
training repetitions and sufficient to justify the additional computational
cost? Classification results, component comparisons and diagnostic
visualisations provide different kinds of evidence for these questions.

This chapter specifies the implemented internal comparison.
Chapter~\ref{chap:external-comparisons} separately reports the completed
single-seed independent implementations of HGTNet and DFAGCN, including
their different preprocessing, optimisers and training budgets. The shared
partitions and evaluation rules do not make these separate recipes a unified
training protocol. Published scores are not treated as measurements made
under the internal protocol; external multi-seed and unified-protocol
comparisons remain open.

\section{Datasets and fixed partitions}

\subsection{NWPU-RESISC45}

Cheng et al.~\cite{cheng2017nwpu} introduced NWPU-RESISC45 as a remote
sensing scene classification benchmark with 45 classes and 31,500 images,
including 700 images per class. The present experiments use a 20\% training
pool and an 80\% held-out test set. The training pool is divided into a
calibration-training subset and a validation subset before final fitting.
Table~\ref{tab:partitions} gives the retained image counts.

The outer partition seed is 42 and the within-pool validation partition seed
is 4242. These fixed partition seeds are distinct from the seeds used for
training repetitions. The manifest construction identified eight groups of
same-class images with identical pixel content. Each such group is kept
within one subset. All original files are retained, so the manifest contains
31,500 files representing 31,492 unique pixel images. This check addresses
the identified exact duplicates; it does not establish the absence of
near-duplicate images or geographic dependence.

\subsection{Aerial Image Dataset}

The Aerial Image Dataset (AID) was introduced by Xia
et al.~\cite{xia2017aid} for aerial scene classification. It contains 10,000
images from 30 scene classes. The implemented partition uses a 50\% training
pool and a 50\% held-out test set. The pool is divided into 4,000 calibration
images and 1,000 validation images. The outer and within-pool partition seeds
are again 42 and 4242. One same-class exact-duplicate group is confined to
one subset; all 10,000 files are retained, with 9,999 unique pixel images.

\subsection{BigEarthNet-19}

Sumbul et al.~\cite{sumbul2019bigearthnet} introduced BigEarthNet as a
multi-label Sentinel-2 archive. The 19-class nomenclature was described by
Sumbul et al.~\cite{sumbul2021bigearthnetmm}. These dataset references do
not define the number of samples retained by this project's available-data
pipeline. The implemented manifests use the existing partition labels in
BigEarthNet v2.0's \texttt{metadata.parquet}~\cite{clasen2024bigearthnetv2}
and retain patches with the required RGB files. They contain 480,038
samples in total, divided as shown in Table~\ref{tab:partitions}; they are not
a newly generated random 70/15/15 split.

Only Sentinel-2 bands B04, B03 and B02 are used, in red, green and blue order.
When separate band files are read, each band value is divided by 4095,
multiplied by 255, clipped to $[0,255]$ and converted to an unsigned 8-bit
image. The reader also supports previously merged RGB images. The inputs
are therefore three-channel RGB representations, rather than full
multispectral or Sentinel-1/Sentinel-2 inputs. Targets use a fixed ordering
of 19 binary land-cover labels.

\begin{table}[tbp]
 \centering
 \caption{Retained sample counts for the implemented internal protocols.
 ``Pool'' is the union of calibration training and validation for the
 single-label datasets. BigEarthNet keeps training and validation separate.}
 \label{tab:partitions}
 \small
 \setlength{\tabcolsep}{4pt}
 \begin{tabular}{lrrrr}
  \toprule
  Dataset & Training & Validation & Pool & Test\\
  \midrule
  NWPU-RESISC45 & 5,040 & 1,260 & 6,300 & 25,200\\
  AID & 4,000 & 1,000 & 5,000 & 5,000\\
  BigEarthNet-19 & 237,871 & 122,342 & --- & 119,825\\
  \bottomrule
 \end{tabular}
\end{table}

\section{Preprocessing and optimisation}

All internal Swin-T and GAViT models receive $224\times224$ RGB inputs.
Their evaluation uses deterministic
resizing and ImageNet channel normalisation, with means
$(0.485,0.456,0.406)$ and standard deviations $(0.229,0.224,0.225)$.
The single-label pipeline uses bilinear resizing with antialiasing. During
single-label training, horizontal and vertical flips each have probability
0.5; brightness, contrast and saturation jitter each use magnitude 0.2.
BigEarthNet training uses the resize and two flips, without colour jitter.
Random training augmentations are disabled during validation and testing.

The internal comparison uses minibatches of 32 and the AdamW optimiser,
which was proposed by Loshchilov and Hutter~\cite{loshchilov2019adamw}.
The initial learning rate is $3\times10^{-4}$ and weight decay is
$10^{-4}$. A cosine schedule uses an annealing horizon $T_{\max}=30$ epochs and a
minimum learning rate of zero. The single-label pipeline uses AdamW moment
coefficients $(0.9,0.999)$ and numerical stabiliser $10^{-8}$. It applies
cross-entropy without label smoothing. BigEarthNet uses binary
cross-entropy with logits. The backbone is fine-tuned together with the
classification and graph modules.

The recorded principal internal runs use 32-bit floating-point computation without
automatic mixed precision (AMP). TensorFloat-32 (TF32) settings are recorded
separately: the single-label pipeline disables TF32 for matrix multiplication
and cuDNN, whereas the documented BigEarthNet soft-seed42 environment disables
matrix-multiplication TF32 and retains cuDNN TF32. These settings describe the
respective pipelines; they do not establish that every historical run used
an identical software or hardware environment.

\section{Model selection and final fitting}

\subsection{Single-label calibration and refitting}

For AID and NWPU-RESISC45, each internal Swin-T or GAViT model first trains for at most 30 calibration
epochs on the internal training subset. Validation overall accuracy selects
the epoch count. The checkpoint is updated only on a strict improvement,
so the earliest occurrence is retained when validation accuracy ties.

After selection, the model is reinitialised from the common pretrained
backbone checkpoint and fitted on the full training pool for the selected
number of epochs. This refit is not continued training from the calibration
checkpoint. The cosine schedule retains its annealing horizon $T_{\max}=30$ epochs;
it is not rescaled to the selected refit duration. The final refit checkpoint
is evaluated on the held-out test set. Thus, a calibration validation score
and a refit test score can come from different trained weights. A reversal
between their rankings alone does not demonstrate overfitting of a single
checkpoint.

\subsection{Multi-label validation-based checkpoint selection}

BigEarthNet uses its training partition for optimisation and its validation
partition for checkpoint selection. The planned training budget is 30 epochs,
and the best checkpoint is updated when validation mean average precision
(mAP) strictly improves. The selected checkpoint is then evaluated on the
fixed test partition. Training and validation are not merged for refitting.
The achieved final epoch and exit status are reported per run, since some
historical runs ended before the planned budget.

The F1 decision threshold is fixed at 0.5. Neither checkpoint choice nor
threshold choice uses the test labels. A higher test F1 score does not
replace mAP as the preselected primary multi-label metric.

\section{Classification metrics}

\subsection{Single-label metrics}

Let $M$ be the number of test images, and let $y_i$ and $\hat{y}_i$ be the
true and predicted class of image $i$. The predicted class is the largest-logit
class. Overall accuracy (OA) is
\begin{equation}
 \operatorname{OA}=\frac{1}{M}\sum_{i=1}^{M}
                       \mathbf{1}[\hat{y}_i=y_i].
 \label{eq:oa}
\end{equation}
For class $c$, let $TP_c$, $FP_c$ and $FN_c$ denote true positives, false
positives and false negatives, respectively. The per-class and macro F1
scores are
\begin{equation}
 F1_c=\frac{2TP_c}{2TP_c+FP_c+FN_c},
 \qquad
 F1_{\mathrm{macro}}=\frac{1}{C}\sum_{c=1}^{C}F1_c.
 \label{eq:macro-f1}
\end{equation}
A zero denominator gives an F1 score of zero in the implemented metric.
Per-class accuracy is $TP_c/(TP_c+FN_c)$ when the class has test support.
The confusion matrix uses true classes as rows and predicted classes as
columns. Scores are displayed as percentages in result tables.

\subsection{Multi-label metrics}

Sigmoid probabilities rank the test images independently for each label.
Let $AP_c$ denote class $c$'s average precision, calculated from that ranking
using the recorded evaluation implementation. For $C_+$ classes with at
least one positive test example, mAP is
\begin{equation}
 \operatorname{mAP}=\frac{1}{C_+}
                 \sum_{c:\,TP_c+FN_c>0}AP_c.
 \label{eq:map}
\end{equation}
The intended full BigEarthNet-19 evaluation covers all 19 labels; absent
classes in small engineering subsets are not grounds for treating subset
mAP as a benchmark result. The class-wise AP vector is retained alongside
the aggregate score.

Binary predictions use $\hat{y}_{ic}=\mathbf{1}[p_{ic}\geq0.5]$, where
$p_{ic}$ is image $i$'s sigmoid probability for class $c$. Macro-F1 uses
Equation~\eqref{eq:macro-f1}. Micro-F1 aggregates counts over all labels:
\begin{equation}
 F1_{\mathrm{micro}}=
 \frac{2\sum_c TP_c}
 {2\sum_c TP_c+\sum_c FP_c+\sum_c FN_c}.
 \label{eq:micro-f1}
\end{equation}
Average precision measures ranking, whereas F1 also depends on the fixed
decision threshold. Their changes can therefore have different directions.

\section{Internal comparisons and repetition statistics}

The principal architecture comparison contains the shared Swin-T baseline,
fixed-partition GAViT, and local-prototype soft-assignment GAViT. Within a
dataset, comparisons retain the same recorded data partitions, label order,
input representation and model-selection rule. The soft variant changes
both region aggregation and the feedback assignment. Accordingly, its
comparison with the fixed variant does not separately isolate the effects
of membership softness, the projection parameters and the feedback map.

Let $v_{ms}$ be the test metric of method $m$ in training repetition $s$, and
let $R$ be the number of eligible repetitions in a stated comparison group.
The reported mean $\bar{v}_m$ and sample standard deviation $s_m$ are
\begin{equation}
 \bar{v}_m=\frac{1}{R}\sum_{s=1}^{R}v_{ms},
 \qquad
 s_m=\sqrt{\frac{1}{R-1}\sum_{s=1}^{R}
                    (v_{ms}-\bar{v}_m)^2},\quad R>1.
 \label{eq:repeat-statistics}
\end{equation}
No standard deviation is assigned to a single run. A paired difference is
computed within the same seed before its mean and standard deviation are
calculated. Differences between percentage scores are expressed in
percentage points (pp).

Each result table states its actual seed set and completion criteria.
Historical runs with different code identities, devices, training durations
or evidence completeness are disclosed separately. They are not relabelled
as fully controlled repetitions. All prespecified runs, including negative
results and interrupted runs, remain in the experiment record. A subset
prepared for a progress report does not replace the complete experiment
record or justify discarding unfavourable runs.

Training repetitions share the frozen test set. They measure variation
conditional on that partition and do not create additional independent test
images or reproduce a repeated-random-split protocol. Means and standard
deviations alone do not establish statistical significance or equivalence.
No significance-test outcome is asserted in this partial draft; the final
test specification must match the completed paired comparisons and their
available sample size.

\section{Efficiency and qualitative analysis}

Efficiency reporting separates parameter count, model-forward latency,
throughput, GPU memory and total experiment duration. A model-forward
benchmark uses the same device, input shape, batch size, numerical precision,
warm-up and synchronisation procedure for all compared models. The archived
AID forward benchmark uses GPU-resident synthetic inputs, 20 warm-up passes
and five groups of 50 timed passes. It excludes image reading, preprocessing
and host-to-device transfer. Its throughput therefore describes the measured
forward path, rather than the complete dataset pipeline. Repeated timing
groups are not independent training seeds, and runs made on different GPU
models are not direct evidence of a model's relative execution speed.

Qualitative analysis distinguishes region assignment, explicit spatial and
feature edges, within-region pooling weights, and graph-attention weights.
These visual quantities represent different operations. Feature similarity
and graph attention are not interchangeable, and neither alone proves a
human-semantic relationship or a causal source of classification improvement.
Anchor locations are distinguished from region centroids. Successes,
introduced errors and shared failures are retained within each declared
selection set. The visualised split and checkpoint identity accompany the
figures, since calibration examples cannot substitute for refit test
performance.

\section{Reproducibility and reporting boundaries}

Run records bind the data manifest, pretrained weight identity, model
configuration, source revision, training seed, numerical settings and
checkpoint to their logs and evaluation outputs. Best validation epoch and
last achieved epoch are recorded separately. A model option in the current
source code is not, by itself, evidence that the corresponding formal
experiment was completed.

Earlier NWPU experiments used a different data split, and earlier graph
variants used different connectivity or integration rules. Those results
are not pooled with the present 20\%-training-pool refit protocol. Likewise,
the historical $3\times3$ BigEarthNet experiment uses a separately frozen
configuration; it does not change the $K=16$ definition of the principal
model in Chapter~\ref{chap:method}. The supplementary external-method runs
in Chapter~\ref{chap:external-comparisons} use their own disclosed
implementation and recipe identities. Their completion does not establish
official reproduction or complete the external multi-seed requirement.
Final eligible seed groups, inferential statistics and complete cross-model
efficiency tables must be resolved against their own evidence before the
thesis makes a final comparative claim.
